\subsection{Why Non-Injectivity Is Not an Ad Hoc Assumption}
  \label{subsec:non-injectivity-not-ad-hoc}

  Non-injectivity of the projection $\Pi:\Omega\to\mathcal{O}$ is sometimes perceived as
  an auxiliary modelling choice introduced to generate desired phenomenology.
  This perception is incorrect; non-injectivity is a structural necessity within any
  relational projection framework that aims to support emergent gauge structure.

  The argument is the following.
  If $\Pi$ were injective and compatible with relational locality, then parallel
  transport of effective states along any closed loop in $\mathcal{O}$ would lift
  uniquely to a closed loop in $\Omega$.
  No two distinct relational configurations would be identified, so the holonomy group
  of the induced fiber bundle would reduce to the identity.
  Trivial holonomy precludes any non-trivial fiber re-identification symmetry, and
  therefore any emergent gauge degree of freedom (Appendix~A.9).
  Injectivity is thus structurally incompatible with the emergence of internal symmetries
  such as the $U(1)$ fiber associated with electromagnetism.

  Non-injectivity is equally compelled from the effective side.
  When the effective description $\mathcal{O}$ reaches a structural saturation --- in the
  sense that it cannot encode all relational distinctions present in $\Omega$ --- multiple
  configurations $\omega\in\Omega$ become indistinguishable under $\Pi$.
  This loss of injectivity is structural rather than epistemic: it reflects an objective
  limitation of the effective representation, not an incomplete specification of underlying
  variables.
  The resulting multiplicity of pre-images $\Pi^{-1}(o)$ for a given $o\in\mathcal{O}$
  defines a structural projection entropy
  \begin{equation}
    S_\Pi \;\equiv\; -\sum_{o\in\mathcal{O}} \nu(o)\,\log\nu(o),
    \qquad \nu(o) = \mu\!\left(\Pi^{-1}(o)\right),
  \end{equation}
  which quantifies the degree of non-injectivity independently of any notion of epistemic
  ignorance.

  In summary, non-injectivity is the unique structural condition that simultaneously
  allows emergent gauge structure (through non-trivial holonomy), finite transport
  capacity (through bounded effective flux), and projective Bell-violation (through
  loss of probabilistic factorization).
  It is a consequence of the relational framework, not an input to it.
